\begin{theo}\label{thm-PP}
Fix $g$, $\ell(\mu)$ and $\ell(\nu)$. Consider the twisted double Hurwitz numbers as a function 
$$\tilde{h}_g(\mu,\nu):\Big\{(\mu,\nu) \mid \sum \mu_i=\sum\nu_j\Big\} \longrightarrow\mathbb{Q}.$$ 
Then the twisted double Hurwitz numbers $\tilde{h}_g(\mu,\nu)$ are piecewise polynomial in the entries $\mu_i$, $\nu_i$. 
The chambers of polynomiality are given by the hyperplane arrangement $\{\sum_{i\in I}\mu_i=\sum_{j\in J}\nu_j\}$, where $I$ and $J$ are subsets of $[\ell(\mu)]$ resp.\ $[\ell(\nu)]$ of size at least~$1$ and at most $\ell(\mu)-1$ resp.\ $\ell(\nu)-1$.
In a chamber, the polynomial $\tilde{h}_g(\mu,\nu)$ is of degree $\ell(\mu)+\ell(\nu)-1+2g$. 
\end{theo}

Here, since we treat $\mu$ and $\nu$ as variables, we assume implicitly that $\Aut(\mu)$ and $\Aut(\nu)$ are trivial. If one inserts special values $\mu,\nu$ for which this is not the case, one has to divide the piecewise polynomial expression by $\Aut(\mu)\cdot \Aut(\nu)$ to pass to the twisted Hurwitz number.

\begin{proof}
The proof builds on the proof of the piecewise polynomiality for double Hurwitz numbers obtained by tropical methods in Theorem 3.1 and 3.9 in~\cite{CJM11}.

Let $\Gamma$ be a twisted monodromy graph of type $(g,\mu,\nu)$. Here, we treat the entries of $\mu$ and $\nu$ as variables, but $g$, $\ell(\mu)$ and $\ell(\nu)$ are fixed numbers. Accordingly, the weights of the edges of $\Gamma$ are also variable and depend on the $\mu_i$ and $\nu_j$. The number $o(\Gamma)$ of vertex orderings compatible with the edge directions, the factor $\frac{1}{|\Aut(\pi)|}$ and the factor $2^{b}$ (where $b$ is the number of branch points) showing up in the multiplicity with which $\Gamma$ contributes to $\tilde{h}_g(\mu,\nu)$ in \cref{cor-monodromygraphs} are all numbers. Thus it remains to show the piecewise polynomiality of the factors $\prod_V (\omega_V-1)$ and $\prod_e\omega(e)$. Recall that the first factor is the product over all $4$-valent vertices $V$ and $\omega_V$ denotes the weight of its adjacent edges, and the second factor the product over all inner edges $e$ and $\omega(e)$ denotes their weight.

Assume $\Gamma$ has $c$ $4$-valent vertices. By an Euler-characteristics computation, $\Gamma$ has $$2(\ell(\mu)+\ell(\nu))+3g-3-c$$ bounded edges and $$ 2(\ell(\mu)+\ell(\nu))+2g-2-2c$$ $3$-valent vertices. The quotient $\Gamma/\iota$ thus has $$\ell(\mu)+\ell(\nu)+\frac{3g-3-c}{2}$$ bounded edges and $$\ell(\mu)+\ell(\nu)+g-c-1$$ $3$-valent vertices. In addition, it has $c$ $2$-valent vertices. We can temporarily remove those and merge the adjacent edges, arriving at a $3$-valent graph with $\ell(\mu)+\ell(\nu)$ ends and $\ell(\mu)+\ell(\nu)+\frac{3g-3-c}{2}-c$ edges. Again by an Euler characteristic computation, for the genus $g'$ of $\Gamma/\iota$ we obtain
$\ell(\mu)+\ell(\nu)+3g'-3$ bounded edges. Equating the two expressions above, we derive at $g'=\frac{1}{2}\cdot (g-c+1)$.

We pick  $g'$ edges in $\Gamma/\iota$ whose removal produces a connected tree. We denote the weights of these edges by variables $i_1,\ldots,i_{g'}$. Then by the balancing condition, the weights of the remaining edges are fixed and a homogeneous linear polynomial in the $\mu_i$, $\nu_j$ and the variables $i_k$. The condition that the edge weights have to be positive provides summation bounds for the variables $i_k$ forming a bounded polyhedron. We sum the product over the edge weights and the factors $(\omega_V-1)$ for the $4$-valent vertices over the lattice points of this polyhedron. We have $\ell(\mu)+\ell(\nu)+\frac{3g-3-c}{2}$ factors for the inner edges, and $c$ for the $4$-valent vertices. Using the Faulhaber formulae for summation of powers of the variables $i_k$, we increase the degree by one with each summation.
Thus, we obtain a polynomial of degree 

\begin{align*}
&\ell(\mu)+\ell(\nu)+\frac{3g-3-c}{2}+c+g'
\\&= \ell(\mu)+\ell(\nu)+\frac{3g-3-c}{2}+c+\frac{1}{2}(g-c+1)
\\&= \ell(\mu)+\ell(\nu)+2g-1.
\end{align*}

As we obtain a polynomial contribution of this degree for every twisted monodromy graph, we deduce that $\tilde{h}_g(\mu,\nu)$ is a polynomial of degree $\ell(\mu)+\ell(\nu)+2g-1$. 

The piecewise polynomial structure and the walls for the chambers of polynomiality arise as in Theorem 3.9 of~\cite{CJM11}, which we state below in~\cref{thm-walls}.
\end{proof}
